\documentclass[11pt]{article}
\usepackage[margin=1in]{geometry}
\usepackage{booktabs}
\usepackage{longtable}
\usepackage{hyperref}
\usepackage{enumitem}

\title{Mango V4 Execution Queue vs MagicBlock Ephemeral Rollups}
\date{March 2026}

\begin{document}
\maketitle

\section{Architecture Comparison}

\begin{longtable}{p{3cm}p{5.5cm}p{5.5cm}}
\toprule
\textbf{Dimension} & \textbf{Mango V4 (Execution Queue)} & \textbf{MagicBlock Ephemeral Rollups} \\
\midrule
Core Model & On-chain ring buffer queue + off-chain deterministic state engine & Off-chain SVM validator + on-chain delegation program \\
\midrule
Execution Location & \textbf{On-chain} (Solana L1 executes every order) & \textbf{Off-chain} (ephemeral SVM processes txs, settles batches to L1) \\
\midrule
State Authority & Solana L1 is always authoritative; harness is optimistic projection & State is cloned off-chain; L1 account is locked until undelegation \\
\midrule
Sequencing & CTM signer assigns monotonic sequence; FIFO enforced on-chain & Single operator/sequencer inside ER; no on-chain ordering guarantee \\
\midrule
Settlement & Implicit --- every order executes as a Solana tx & Explicit --- periodic CommitState + challenge window + Finalize \\
\bottomrule
\end{longtable}

\section{Key Tradeoffs}

\subsection{Where MagicBlock Wins}

\begin{itemize}
  \item \textbf{Latency}: 10--50ms per tx (no consensus overhead) vs 400ms--1.2s (Solana slot time)
  \item \textbf{Throughput}: Unbounded horizontal scaling (spin up more ERs) vs 1,024-item ring buffer bottleneck
  \item \textbf{Compute costs}: Near-zero fees inside ER; Mango pays full Solana CU per order (\textasciitilde30--50K CU each)
  \item \textbf{No stall problem}: Execution isn't bottlenecked by L1 compute budgets (STALL-1/STALL-4 don't exist)
\end{itemize}

\subsection{Where Mango V4 Wins}

\begin{itemize}
  \item \textbf{No new trust assumptions}: Crankers are fully untrusted; MagicBlock requires trusting the ER operator for ordering/censorship resistance during the session
  \item \textbf{No fraud-proof complexity}: Settlement is atomic per-tx on L1 --- no challenge windows, no dispute resolution, no data availability assumptions
  \item \textbf{Full L1 composability during execution}: Every order is a native Solana tx composable with all of DeFi; MagicBlock locks accounts on L1 while delegated
  \item \textbf{Deterministic auditability}: Dual-signature scheme (user intent + CTM envelope) creates a cryptographic audit trail fully verifiable on-chain; MagicBlock's off-chain execution is opaque until commitment
  \item \textbf{Simpler recovery}: Divergence = re-read L1 state. MagicBlock recovery involves fraud proofs, forced undelegation, and potential state reversion
\end{itemize}

\section{Security Model Comparison}

\begin{longtable}{p{3cm}p{5.5cm}p{5.5cm}}
\toprule
\textbf{Concern} & \textbf{Mango V4} & \textbf{MagicBlock} \\
\midrule
Sequencer censorship & CTM signer can censor; mitigated by direct-submit bypass (C-4 fix, 10-slot delay) & ER operator can censor; mitigated by time-bound delegation auto-expiry \\
\midrule
Execution integrity & Verified per-tx on L1 (payload hash + accounts hash) & Optimistic --- trust operator, then fraud-proof window \\
\midrule
Data availability & All data on-chain (Solana is the DA layer) & Operator must publish data; if withheld, fraud proofs can't be constructed \\
\midrule
Account safety & Accounts never leave L1 ownership & Accounts change owner to Delegation Program; locked on L1 \\
\midrule
MEV protection & Strict FIFO + cryptographic binding prevents reordering & Operator has full reordering power within ER \\
\bottomrule
\end{longtable}

\section{Fundamental Architectural Difference}

Mango V4 is a \textbf{sequenced execution queue on L1} --- accepting L1 latency/throughput constraints but gaining zero new trust assumptions. Every order is a real Solana transaction.

MagicBlock is an \textbf{off-chain execution environment} --- breaking free of L1 constraints but introducing operator trust, fraud proofs, and state locking complexity.

\section{Potential Hybrid Approach}

The two architectures are not mutually exclusive:

\begin{enumerate}
  \item \textbf{Use MagicBlock ERs as the execution engine} --- delegate the execution queue + orderbook accounts to an ER, process orders at 10--50ms, then settle batches back
  \item \textbf{Keep CTM sequencing layer} --- MagicBlock's single-operator model lacks cryptographic ordering guarantees; CTM envelope + FIFO could run \emph{inside} the ER
  \item \textbf{Selective delegation} --- use ERs only for high-frequency markets, keep lower-volume markets on the current L1 queue
\end{enumerate}

This would yield MagicBlock's latency (\textasciitilde50ms) with Mango's security properties (cryptographic sequencing, permissionless verification, no MEV reordering).

\section{Conclusion}

Mango V4 prioritizes \textbf{trustlessness and auditability} over raw speed. MagicBlock prioritizes \textbf{performance} at the cost of new trust assumptions. For a DeFi protocol handling real money, the conservative approach has merit --- but the latency gap is significant for competitive market making.

\end{document}
